\paragraph{Task vectors and per-task geometry.} Fix integers
$T\geq 1$, $d\geq 2$, $1\leq r\leq d$, and a radius $B > 0$.
Each of the $T$ tasks is specified by a pair $(\tau_t, H_t)\in
\R^d\times\R^{d\times d}$ with
\begin{equation}\label{eq:admissibility}
  H_t \succeq 0,\quad \rank(H_t) \leq r,\quad \tau_t \in V_t :=
  \range(H_t),\quad \tau_t^\top H_t \tau_t \leq B^2.
\end{equation}
Per-task loss is the quadratic $f_t(w) := \tfrac12 (w - \tau_t)^\top
H_t (w - \tau_t)$, which attains zero at any $w$ with $P_{V_t} w
= \tau_t$. We write $\boldsymbol\tau := (\tau_t)_{t=1}^T$ and
$\boldsymbol H := (H_t)_{t=1}^T$ for the tuple data, and call any
$(\boldsymbol\tau, \boldsymbol H)$ satisfying
\eqref{eq:admissibility} \emph{admissible}.

The specialization $H_t = P_{V_t}$ (the orthogonal projector onto
$V_t$) recovers the weight-disentangled LoRA setting of
Ortiz-Jimenez et al.~\cite{ortizjimenez2023disentanglement}; the
general $H_t\succeq 0$ case subsumes Fisher-metric surrogates for
cross-entropy LLM losses with heterogeneous task-specific
curvature.

\paragraph{Worst-task distortion.} A merged model $w^\star\in
\R^d$ incurs per-task distortion $D_t(w^\star) := (w^\star-
\tau_t)^\top H_t (w^\star - \tau_t)$. We take as figure of merit
the \emph{worst-task} distortion
\begin{equation}\label{eq:Delta}
  \Delta(w^\star; \boldsymbol\tau, \boldsymbol H)
  \;:=\; \max_{t\in[T]}\; D_t(w^\star).
\end{equation}
The choice of $\max$ over $\mathrm{avg}$ matters. A merged model
that performs well on most tasks but catastrophically on one is
not a useful deployment target; classical multi-description
coding~\cite{elgamalcover1982md} treats average distortion and
has a strictly weaker rate-distortion function than the one we
establish (see Lemma~\ref{lem:id}).

\paragraph{Rate-$R$ merging codes.} A rate-$R$ merging code is a
pair $(E, D)$ with encoder $E : (\R^d)^T \to \{1,\ldots,2^R\}$
and decoder $D : \{1,\ldots,2^R\}\to\R^d$. The encoder observes
$\boldsymbol\tau$ (the $H_t$'s may be assumed known to the
encoder as side information; our lower bound holds even if the
encoder is told $\boldsymbol H$ for free). The reconstruction is
$w^\star := D(E(\boldsymbol\tau))$. The \emph{rank-$r$
rate-distortion function} is
\begin{equation}\label{eq:rd-def}
  \mathcal{D}^\star(T, d, r, B, R)
  \;:=\; \inf_{(E, D)}\;\sup_{(\boldsymbol\tau, \boldsymbol H)
    \text{ admissible}}\;
  \E\bigl[\Delta(w^\star;\boldsymbol\tau, \boldsymbol H)\bigr],
\end{equation}
where the expectation is over any randomness the encoder and
decoder share.

With the figure of merit, admissibility constraints, and code
class fixed, the central question of the paper is the following.

\begin{problem}[Rank-$r$ rate-distortion model merging]
\label{prob:merging}
Determine the rank-$r$ rate-distortion function
$\mathcal{D}^\star(T, d, r, B, R)$ defined in
\eqref{eq:rd-def}, and exhibit an explicit encoder/decoder pair
$(E, D)$ attaining it. Concretely:
\begin{enumerate}[leftmargin=*,topsep=2pt,itemsep=0pt]
  \item (\emph{Lower bound}) Find the largest $L(T, d, r, B, R)$
  such that for every rate-$R$ merging code,
  $\sup_{(\boldsymbol\tau,\boldsymbol H)\text{ admissible}}
  \E[\Delta(w^\star;\boldsymbol\tau,\boldsymbol H)]\geq
  L(T, d, r, B, R)$.
  \item (\emph{Upper bound}) Construct a rate-$R$ merging code
  $(E^\star, D^\star)$ achieving distortion at most
  $U(T, d, r, B, R)$ on every admissible tuple.
  \item (\emph{Tightness}) Establish $L = \Theta(U)$.
\end{enumerate}
The remainder of the paper resolves (i)
(Theorem~\ref{thm:general}, \S\ref{sec:lower}), (ii)
(Theorem~\ref{thm:achv}, \S\ref{sec:achv}), and (iii) up to a
universal constant in the floor-zero Stiefel regime.
\end{problem}

\paragraph{Key derived quantities.} We will repeatedly use
\begin{equation}\label{eq:derived}
  \bar H := \tfrac1T\sum_{t=1}^T H_t,\qquad
  \tauH := \bar H^+\cdot \tfrac1T\sum_{t=1}^T H_t\tau_t,\qquad
  \deff := \rank\!\Bigl(\sum_{t=1}^T P_{V_t}\Bigr),
\end{equation}
where $(\cdot)^+$ denotes the Moore--Penrose pseudoinverse. The
\emph{$H_t$-weighted centroid} $\tauH$ is the unique element of
$\range(\bar H) = V_1 + \cdots + V_T$ satisfying $\bar H\tauH =
T^{-1}\sum_t H_t\tau_t$; it reduces to the arithmetic mean
$\bar\tau$ when $H_t = I_d$. The \emph{effective dimension}
$\deff \leq \min(Tr, d)$ is the combinatorial quantity that
controls the rate-distortion function; its meaning is the number
of independent directions the decoder must reconstruct.

\paragraph{Hard distribution $P^\star$.} Our lower bound is
established against a specific distribution $P^\star$ on
admissible tuples. For each task $t\in[T]$, independently:
\begin{enumerate}[leftmargin=*,topsep=2pt,itemsep=2pt]
  \item Sample $U_t\in\mathrm{Stiefel}(d, r)$ uniformly
  (equivalently, take the Gram--Schmidt orthogonalization of a
  $d\times r$ matrix of iid Gaussians); set $V_t := \range(U_t)$.
  \item Fix a positive-definite diagonal $r\times r$ matrix
  $D_t\succ 0$ (the task-$t$ spectrum inside $V_t$); set $H_t :=
  U_t D_t U_t^\top$.
  \item Sample $u_t\sim\mathrm{Unif}(S^{r-1})$ and set
  $v_t := B\, D_t^{-1/2} u_t \in\R^r$, $\tau_t := U_t v_t$.
\end{enumerate}
The scale normalization $\tau_t^\top H_t\tau_t = v_t^\top D_t
v_t = B^2\|u_t\|^2 = B^2$ holds by construction. The $D_t^{-1/2}$
stretch makes $\tau_t$ uniform on the $H_t$-weighted unit
ellipsoid inside $V_t$, which is what makes the second-moment
identity $\E[H_t\tau_t\tau_t^\top H_t | U_t] = (B^2/r) H_t$ hold
\emph{regardless of} $D_t$. This $D_t$-freeness is the feature
that makes Theorem~\ref{thm:general} apply to arbitrary
task-dependent spectra. For the projector special case $H_t =
P_{V_t}$ (i.e.\ $D_t = I_r$), $P^\star$ reduces to the
Stiefel-random LoRA distribution with $\tau_t\sim\mathrm{Unif}
(B\cdot S^{r-1}\cap V_t)$.

\paragraph{Yao's minimax bridge.} The worst-case rate-distortion
function $\mathcal{D}^\star$ in \eqref{eq:rd-def} is at least the
Bayesian rate-distortion function under any prior. In particular,
taking the prior to be $P^\star$:
\begin{equation}\label{eq:yao}
  \mathcal{D}^\star(T,d,r,B,R)
  \;\geq\;
  \inf_{(E,D)}\;\E_{(\boldsymbol\tau,\boldsymbol H)\sim P^\star}
    \bigl[\Delta(w^\star;\boldsymbol\tau,\boldsymbol H)\bigr].
\end{equation}
This is Yao's minimax principle~\cite{yao1977} applied to the
randomized-code / worst-case-input game. All subsequent lower
bounds are proved against the RHS of \eqref{eq:yao}; the LHS
inherits them for free.
